El procedimiento sigue las cinco fases del modelo (Anexo~C), alineadas con OE1--OE4. La
\textbf{adquisición de artefactos} (Cap.~4 §4.3) precede a la normalización: contratos OpenAPI y
material BIAN se identifican, versionan y almacenan antes del \textit{parsing}. Este capítulo
documenta el diseño del procedimiento; su ejecución a escala queda fuera del alcance (§1.6).

\begin{enumerate}
  \item \textbf{Fase 1 --- Diseñar el modelo:} definir entradas, scores S/E/C, fórmula de
        agregación y umbrales de clasificación (Cap.~1 §1.4; Anexo~C).
  \item \textbf{Fase 2 --- Adquirir y normalizar artefactos (OE1):} obtener contratos OpenAPI y
        extractos BIAN (§4.3); parsear y producir representaciones comparables (tabular, grafo o RDF
        según factor elegido).
  \item \textbf{Fase 3 --- Emparejar estructuras (OE2):} calcular correspondencias y
        \StructuralScore{} (E).
  \item \textbf{Fase 4 --- Calcular similitud semántica (OE3):} obtener \SemanticScore{} (S).
  \item \textbf{Fase 5 --- Integrar y clasificar (OE4):} calcular \CoverageScore{} (C),
        \AlignmentScore{} y clasificación por tipologías (Tabla~\ref{tab:tipologias}).
\end{enumerate}

\subsection{Relación objetivos, fases y procedimiento}
\label{sec:objetivos-metodologia}

Esta subsección explicita cómo cada objetivo específico se desarrolla en el procedimiento. La
\Cref{tab:oe-metodologia} vincula OE, fase, acción metodológica y producto verificable.

\begin{table}[H]
  \centering
  \caption{Objetivos específicos vinculados al procedimiento metodológico}
  \label{tab:oe-metodologia}
  \small
  \begin{tabular}{@{}p{0.06\textwidth}p{0.14\textwidth}p{0.22\textwidth}p{0.24\textwidth}p{0.24\textwidth}@{}}
    \toprule
    \textbf{OE} & \textbf{Fase} & \textbf{Acción (verbo)} & \textbf{Entrada} & \textbf{Producto} \\
    \midrule
    --- &
    Fase 1 &
    \textbf{Diseñar} entradas, scores y umbrales &
    Requisitos del §1.4; Anexo C &
    Modelo S/E/C documentado \\
    \textbf{OE1} &
    Fase 2 &
    \textbf{Adquirir} y \textbf{normalizar} contrato y SD &
    OpenAPI + extracto BIAN (§4.3) &
    Representaciones comparables \\
    \textbf{OE2} &
    Fase 3 &
    \textbf{Emparejar} estructuras y calcular E &
    Representaciones OE1 &
    Matriz de correspondencias + E \\
    \textbf{OE3} &
    Fase 4 &
    \textbf{Calcular} similitud semántica S &
    Representaciones OE1 &
    Matriz semántica + S \\
    \textbf{OE4} &
    Fase 5 &
    \textbf{Integrar} S, E, C y clasificar &
    E, S, representaciones OE1 &
    C, \AlignmentScore{}, tipologías \\
    \bottomrule
  \end{tabular}
\end{table}

Cada fila corresponde a un hito del Capítulo~\ref{cap:objetivos} y a la matriz de consistencia
(Anexo~C). Las técnicas concretas por OE se detallan en la \Cref{sec:tecnicas-herramientas}.


La \Cref{fig:proc-estructural} ilustra la correspondencia jerárquica para E; la cadena completa
aparece en la \Cref{fig:cadena-objetivos} (Cap.~2).

\setcounter{figure}{6}
\begin{figure}[H]
  \centering
  \figCapSeven
  \caption{Correspondencia jerárquica para el cálculo de \StructuralScore{} (E). Fuente: elaboración
  propia; ejemplo ilustrativo.}
  \label{fig:proc-estructural}
\end{figure}

\subsection{Protocolo de validación empírica (fuera de alcance)}

Para una validación empírica posterior, el protocolo escrito contemplará: (1)~instrumento de captura
(repositorio de contratos); (2)~infraestructura (\textit{parsing}, almacenamiento);
(3)~condiciones de registro; (4)~tamaño y representatividad de la muestra; (5)~etiquetado de pares
OpenAPI$\leftrightarrow$BIAN (juicio experto o \textit{ground truth} documentado);
(6)~aspectos éticos/legales; (7)~evaluación del desempeño del sistema (precisión vs.\ referencia,
análisis de falsos positivos/negativos en matching).
